\cleardoublepage

% concept03.tex
\section{Comparing and Ordering Numbers}

Math is not just about knowing what numbers are. It is about knowing how they
relate. We compare. We rank. We ask, which is bigger? Which is smaller? Which is
equal? Symbols like \(>\), \(<\), and \(=\) help us answer. But first, we have to
understand the values.

Take \(5.03\) and \(5.3\). They may look close. But \(5.3\) is larger. Now try
comparing fractions: \(\frac{2}{5}\) and \(\frac{3}{7}\). Convert them to
decimals. \(\frac{2}{5}\) is \(0.4\). \(\frac{3}{7}\) is about \(0.428\). So
\(\frac{3}{7}\) is greater. This kind of comparison matters in real life--from
prices to measurements to data.

Ordering numbers builds the same skill. Imagine you are given \(0.2\),
\(\frac{1}{3}\), \(0.25\), and \(0.199\). Which is the greatest? Which is least?
Convert, compare, and you find the order. This process teaches more than math. It
teaches precision, logic, and clarity.

\blockquote{Math is not just about knowing what numbers are. It is about knowing
how they relate.}

\subsection{How do we compare and order numbers?}

In mathematics, it's important not only to understand numbers, but also to compare
them and put them in order. We use symbols like \(>\) (greater than), \(<\) (less
than), and \(=\) (equal to) to show how numbers relate to each other. For example,
if you are given two numbers, you can decide which is larger or smaller, or if
they are the same.

Let's look at a few examples. Suppose you want to compare \(5.03\) and \(5.3\).
Since \(5.03 < 5.3\), we use the less than symbol. Now, what if you have two
fractions, \(\frac{2}{5}\) and \(\frac{3}{7}\)? By converting them to decimals,
you find that \(\frac{2}{5} = 0.4\) and \(\frac{3}{7} \approx 0.428\), so
\(0.4 < 0.428\), which means \(\frac{2}{5} < \frac{3}{7}\).

Ordering numbers is also a useful skill. Imagine you are asked to arrange the
numbers \(0.2,\ \frac{1}{3},\ 0.25,\ 0.199\) from greatest to least. By
converting them all to decimals, you see that \(\frac{1}{3}\) (approximately
\(0.333\)) is the largest, followed by \(0.25\), then \(0.2\), and finally
\(0.199\). Sometimes, you may need to compare a fraction and a decimal, such as
\(\frac{4}{9}\) and \(0.44\). Since \(\frac{4}{9} \approx 0.444\), it is slightly
greater than \(0.44\). These skills are essential for working with numbers in
real-life situations.

\blockquote{Take \(5.03\) and \(5.3\). They may look close. But \(5.3\) is
larger.}


% \cleardoublepage
\subsection{titlehere (English)}

question here

\fullpageimage{images/q1_3_1-eng.jpg}


\newpage
\subsection{titlehere (Korean)}

question here

\fullpageimage{images/q1_3_1-kor.jpg}


\newpage
\subsection{titlehere (Answer)}

question here

% \cleardoublepage
\subsection{Long Question (English)}
Put the following numbers in order from least to greatest: 0.199, $\frac{1}{3}$, 0.25, 0.2

\fullpageimage{images/q1_3_2-eng.jpg}

\newpage
\subsection{Long Question (Korean)}
Put the following numbers in order from least to greatest: 0.199, $\frac{1}{3}$, 0.25, 0.2

\fullpageimage{images/q1_3_2-kor.jpg}

\newpage
\subsection{Long Question (Answer)}
Put the following numbers in order from least to greatest: 0.199, $\frac{1}{3}$, 0.25, 0.2

\textbf{Answer:}\\
Convert to decimals:\\
- 0.199\\
- $\frac{1}{3} \approx 0.333$\\
- 0.25\\
- 0.2\\
\textbf{Ordered:} 0.199 < 0.2 < 0.25 < $\frac{1}{3}$

% \cleardoublepage
\subsection{Challenge Question (English)}
Which is greater: $\frac{4}{9}$ or 0.44? Show your reasoning.

\fullpageimage{images/q1_3_3-eng.jpg}

\newpage
\subsection{Challenge Question (Korean)}
Which is greater: $\frac{4}{9}$ or 0.44? Show your reasoning.

\fullpageimage{images/q1_3_3-kor.jpg}

\newpage
\subsection{Challenge Question (Answer)}
Which is greater: $\frac{4}{9}$ or 0.44? Show your reasoning.

\textbf{Answer:}\\
$\frac{4}{9} \approx 0.444...$ which is slightly greater than 0.44.\\
\textbf{Answer:} $\frac{4}{9}$ is greater.

